\section{Relevant results}\label{sec:results}

\textit{What follows are the concrete findings from the literature that my own system and evaluation will be measured against, or will try to build on / improve on, organised by what each cluster of results is relevant to.}



\subsection{Citation-grounded / attributed generation methods — results relevant to \textit{which generation architecture to adopt}}\label{sec:res-generation}

\begin{itemize}
\item \textbf{RARR} (Google, 2022) — post-hoc: researches evidence via web search for any LLM's claims, edits unsupported text. Model-agnostic, few-shot. arXiv:2210.08726
\item \textbf{ALCE} (Princeton/Gao et al., EMNLP 2023) — benchmark + pipeline for citation generation. ChatGPT: 73.6\%/72.5\% citation recall/precision on ASQA, improving to 84.8\%/81.6\% with reranking; collapsed to $\sim$20\% on harder list-style QAMPARI task. \textquotedblleft{}Even the best models lack complete citation support 50\% of the time\textquotedblright{} on ELI5. \textbf{Relevance:} this is the precision/recall definition my evaluation harness implements structurally; the ASQA-vs-QAMPARI gap is a warning that citation quality is highly task-dependent, so InnoBeweegLab's own task shape must be tested directly rather than assumed from a generic benchmark. arXiv:2305.14627
\item \textbf{Self-RAG} (Asai et al., UW/AI2, ICLR 2024) — trains model to decide when to retrieve + emit self-critique \textquotedblleft{}reflection tokens\textquotedblright{} rather than always retrieving fixed passages. 13B model: 70.3\%/71.3\% citation precision/recall on ASQA vs. naive retrieval-augmented Llama2-13B at only 19.8\%/36.1\%. \textbf{Relevance: architecture/training for citation matters far more than bolting retrieval onto a normal LLM} — directly informs SQ3's architecture comparison. arXiv:2310.11511
\item \textbf{FLARE} (CMU, EMNLP 2023) — iterative, triggers retrieval only when the model is about to generate a low-confidence token. arXiv:2305.06983
\item \textbf{AIS framework} (Rashkin et al., Google, 2021) — formalized \textquotedblleft{}Attributable to Identified Sources\textquotedblright{} (AIS), a framework for judging whether generated text is attributable to given source material; published in \textit{Computational Linguistics} (MIT Press, 2023). The related \textbf{Attributed QA} benchmark (Bohnet et al., Google, 2022) applies this framework to open-domain QA with retrieval-augmented systems. arXiv:2112.12870 (AIS) / arXiv:2212.08037 (Attributed QA)
\item \textbf{FacTool} (Chern et al., 2023) — general claim-extraction $\rightarrow$ verification pipeline across QA/code/math/science. arXiv:2307.13528
\item \textbf{Chain-of-Verification (CoVe)} (Meta, ACL Findings 2024) — draft $\rightarrow$ generate independent verification questions $\rightarrow$ answer blind $\rightarrow$ revise. arXiv:2309.11495
\end{itemize}

\subsection{Automated factuality/faithfulness evaluation — results relevant to \textit{how well my own evaluation method can be trusted}}\label{sec:res-factuality}

\begin{itemize}
\item \textbf{FActScore} (UW/Meta, EMNLP 2023) — decomposes text into atomic facts, checks each against a KB. ChatGPT scored only \textbf{58\%} on biography generation. Automated estimator had <2\% error vs. human eval, at a fraction of the cost (\$26K human-eval cost avoided). arXiv:2305.14251 — \textit{re-verified directly against the arXiv abstract on 2026-10-07; this is also the paper the evaluation harness's \texttt{faithfulness} function (\texttt{src/evaluation/metrics.py}) is modelled on.}
\item \textbf{SummaC} (Laban et al., TACL 2022) — repurposes NLI models, sentence-pair aggregation. 74.4\% balanced accuracy, +5pts over prior methods. arXiv:2111.09525
\item \textbf{TRUE benchmark} (Google, NAACL 2022) — meta-evaluates factual-consistency metrics; NLI-based and QA-based approaches are strongest and complementary. arXiv:2204.04991
\item \textbf{RAGAS} (Es et al., EACL 2024) — open-source, industry-standard RAG evaluation. \textbf{Faithfulness} = claims supported by retrieved context / total claims. Also Context Precision/Recall auditing the retrieval step itself. Already used as evaluation framework in the Chalmers/Volvo thesis reviewed earlier (Section~\ref{sec:res-prior}). arXiv:2309.15217
\item \textbf{SAFE + LongFact} (Google DeepMind/Stanford, 2024) — atomic-decomposition grounded against live search. Agreed with human raters 72\% of the time; on disagreements, SAFE was correct 76\% of the time; >20x cheaper than human annotation. \textbf{Relevance:} strong evidence automated fact-checking can match/exceed human-rater reliability on a budget — directly supports using an automated harness as the primary evaluation method, with human spot-checks rather than full human re-grading. arXiv:2403.18802
\item \textbf{AttributionBench} (2024) — benchmarks whether LLMs can reliably judge if a citation supports a claim. Ceiling $\sim$80\% macro-F1 even for fine-tuned GPT-3.5; bigger model (GPT-4) didn't help much. See Section~\ref{sec:gap}. arXiv:2402.15089
\end{itemize}

\subsection{Production systems and audits — results relevant to \textit{what traceability failure actually looks like in deployed systems}}\label{sec:res-production}

\begin{itemize}
\item \textbf{Anthropic Citations API} (2025) — chunks source docs, attaches every response to source chunks. Case study (Endex): source hallucinations 10\%$\rightarrow$0\%, +20\% references/response. Vendor-reported.
\item \textbf{CJR/Tow Center audit} (2024/25) — 8 generative search engines, 1,600 queries. Perplexity 37\% incorrect; ChatGPT Search >60\% incorrect; Grok-3 94\% error rate. \textbf{Paid tiers were more confidently wrong than free tiers.}
\item \textbf{AlgorithmWatch Bing Chat study} (2023) — $\sim$30\% factual errors in election Q\&A; fabricated claims attributed to real, credible outlets — \textquotedblleft{}illusion of validity\textquotedblright{} via false attribution.
\item \textbf{Stanford RegLab legal-AI study} (2025, preregistered) — Lexis+ AI and Ask Practical Law AI both hallucinate in >17\% of responses despite \textquotedblleft{}hallucination-free\textquotedblright{} marketing. Source of the misgrounding/fabrication distinction used throughout this review (Section~\ref{sec:definitions}).
\item \textbf{Vals Legal AI Report} (2025) — Harvey/CoCounsel vs. practicing lawyers; Harvey 94.8\% on document Q\&A. Citation/traceability performance is highly task-dependent, not one global capability.
\item \textbf{Glean} — enterprise RAG, permission-aware, cites the Stanford legal numbers as rationale for grounding.
\item \textbf{\textquotedblleft{}Cited but Not Verified\textquotedblright{}} (2026) — 14 deep-research agents, 130 queries: link-validity 94\%+, topical-relevance 80\%+, factual-accuracy only \textbf{39–77\%}. Fact-checking accuracy degrades $\sim$42\% as agents take more search steps. \textbf{More citations per response correlated with LOWER factual accuracy.} Single most directly on-topic recent study for this thesis. arXiv:2605.06635
\end{itemize}

\subsection{Structured-extraction \& multi-agent report pipelines — results relevant to \textit{whether InnoBeweegLab's use case needs a multi-agent architecture}}\label{sec:res-pipelines}

\begin{itemize}
\item \textbf{Auto-TA} (2025) — closest analog to InnoBeweegLab's use case. Multi-agent qualitative-coding pipeline with a Quote ID on every code/theme. Credibility rose from 82.13\% baseline to \textbf{98.41\%} with a specialized persona + grounding-checker agent. arXiv:2506.23998
\item \textbf{SFT-TA} (2025) — fine-tuned agents underperform GPT-4o alone individually, but outperform both fine-tuned-single-agent and GPT-4o baseline once embedded in a multi-agent system. arXiv:2509.17167
\item \textbf{PublicAgent} (Montazeri, Feng \& Sha, 2026) — a multi-agent architecture for open-data analysis (intent clarification, data discovery, data analysis, report generation agents) tested via ablation across 8 LLMs and 50 benchmark queries: removing any single agent cost the full system 72–100\% of pairwise wins against the ablated version, and different agents failed in distinct, non-redundant ways — discovery/analysis-agent removal caused catastrophic failure (empty or nonsensical output), while report/intent-agent removal caused milder quality degradation. Specialisation improved results independent of the base model's own strength. \textbf{Relevance:} a finer-grained empirical picture for SQ3 than a single pass/fail verdict — discovery/analysis-type agents (handling data, closest to InnoBeweegLab's SourceRecord/ingestion layer) were consistently valuable, while report/intent-type agents were more model-dependent, echoing the caution the Chalmers/Volvo result (Section~\ref{sec:res-prior}) raises specifically for report-synthesis agents. doi.org/10.1016/j.dte.2026.100136
\item \textbf{GATOS Workflow} (Katz, Coloyan Fleming \& Main, 2024) — open-source, locally-runnable thematic-analysis pipeline (Llama-3.1-70B/8B, WizardLM2-8x22B, Mistral-22B/12B-nemo), explicitly avoiding cloud APIs for sensitive data. Validated only on three synthetic case-study datasets: fewer than 5\% of sub-themes had no clear match, under 2\% had no close match at all. arXiv:2410.03721
\item \textbf{Nature HSSC 2026} (Katz, Coloyan Fleming \& Main) — the peer-reviewed version of the GATOS workflow above, published in \textit{Humanities and Social Sciences Communications} 13, article 209.
\item \textbf{GATOS applied to real data} (Kim, Ke, Zhang \& Barrett, EDM 2025) — a separate paper applying the Katz et al. GATOS workflow to real (not synthetic) qualitative data, using DeepSeek + mxbai-embed-large-v1 embeddings, on a neurodiversity-support study (interviews with healthcare providers, teachers, caregivers). This narrows, but doesn't close, the real-data validation gap named below: it still used hosted inference, not InnoBeweegLab's 16GB-RAM, zero-budget, fully local constraint.
\item \textbf{PaperTrail} (CHI 2026) — bi-directional claim decomposition (both the LLM answer AND source docs), flags \textquotedblleft{}supported / unsupported / omitted.\textquotedblright{} arXiv:2602.21045
\end{itemize}

\subsection{Human-in-the-loop verification interface patterns — results relevant to \textit{designing the Workshop-phase review interface}}\label{sec:res-hitl}

\begin{itemize}
\item \textbf{PaperTrail} (CHI 2026) — 26 scholars, within-subjects study. Structured claim-evidence UI \textbf{lowered trust appropriately but did NOT change verification behavior} — reviewers still mostly accepted LLM edits.
\item \textbf{\textquotedblleft{}Hide or Highlight\textquotedblright{}} (AIES 2025/26, N=148) — hiding/vaguing uncertain content produced \textit{higher and more appropriate} trust than highlighting it (appropriate compliance: 90\% Ambiguity / 86\% Opaque vs. only 66\% Transparent highlighting).
\item \textbf{AutoUAT / Test Flow} (Ferreira, Viegas, Faria \& Lima, 2025 industrial case study) — a two-step LLM pipeline generating acceptance-test scenarios (95\% rated helpful by product owners, across 166 real usage instances) and executable test scripts (60\% usable as generated, rising to 92\% with minor fixes) from user stories, deployed in an automotive-industry workflow with mandatory human review before any generated test merges. Different domain from report generation, but the same pattern as PaperTrail/\textquotedblleft{}Hide or Highlight\textquotedblright{}: generated output is treated as a draft a human must approve, not a final artefact — directly supporting UC4's human-in-the-loop requirement. The tool's output also once surfaced a real application defect the generated test scenario had flagged but the human team hadn't noticed, a concrete example of traceable AI output adding value beyond pure time-saving. arXiv:2504.07244
\item Broader CSCW/CHI literature: showing confidence scores doesn't reliably improve human error-catching, and can increase over-reliance.
\end{itemize}

\subsection{Prior comparable work — results relevant to \textit{hardware and architecture feasibility specifically}}\label{sec:res-prior}

\begin{itemize}
\item The \textbf{Chalmers/Volvo master's thesis} (Nihlstrand \& Lundberg, 2025, multi-agent LLM system for automated report generation from unstructured logbooks) used Phi-4 (14B), LLaMA-3.3-70B and GPT-4o, all via hosted inference, and found its multi-agent system scored \textit{worse overall} than a naive single-pass baseline (3.60/7 vs 3.78/7), only beating it on factuality specifically. It also found RAG/retrieval grounding — not raw model size — was the biggest lever on factuality in its own ablation study, which this review's own findings (Section~\ref{sec:res-generation}, Self-RAG) independently corroborate. It did not validate laptop-class local deployment at all; a 70B model is far outside InnoBeweegLab's 16GB-RAM laptop.
\end{itemize}

\subsection{Systematic evaluation methodology as its own field of study — results relevant to \textit{designing a trustworthy evaluation protocol in general, not just assembling a project-specific checklist}}\label{sec:res-evalmethod}

Sections~\ref{sec:res-generation}--\ref{sec:res-prior} above answer \textquotedblleft{}which metrics and generation methods exist.\textquotedblright{} This subsection asks a different, prior question, raised directly by feedback that the methodology in Section~\ref{sec:m-reference} needed to read as literature-grounded research rather than an invented one-off: how does the literature say an evaluation \textit{protocol itself} should be designed? Four sub-threads of this field are directly relevant.



\textbf{Evaluation as a three-part design problem, not a single metric choice.} Chang et al.'s \textquotedblleft{}A Survey on Evaluation of Large Language Models\textquotedblright{} (arXiv:2307.03109, ACM TIST 2024) frames systematic evaluation of any generative system around three separable questions — \textit{what} to evaluate (which capabilities/tasks matter for the use case), \textit{where} to evaluate (which dataset/benchmark/gold set actually represents that use case), and \textit{how} to evaluate (automatic metric, human judgment, or LLM-as-judge, each with different cost/reliability trade-offs). \textbf{Relevance:} this is the structure Section~\ref{sec:m-reference} is implicitly following (what = traceability/faithfulness/content coverage; where = the frozen gold set plus the real Euterpepark report and a naive baseline; how = structural metrics + human rubric), but had not previously been named as a methodology choice in its own right.



\textbf{LLM-as-judge is its own well-studied sub-field, with documented, specific failure modes — not a generic \textquotedblleft{}AI checks AI\textquotedblright{} shortcut.} Gu et al.'s \textquotedblleft{}A Survey on LLM-as-a-Judge\textquotedblright{} (arXiv:2411.15594) proposes a formal reliability framework built on three axes — agreement with human judgment, susceptibility to bias, and adversarial robustness — and separates \textit{task-agnostic} biases (diversity, cultural, self-enhancement bias) from \textit{judgment-specific} ones (position, length/verbosity, style, concreteness bias). Zheng et al.'s earlier \textquotedblleft{}Judging LLM-as-a-Judge with MT-Bench and Chatbot Arena\textquotedblright{} (arXiv:2306.05685) is the empirical paper this taxonomy builds on, and measured these directly in GPT-4-as-judge setups: \textbf{position bias} (the judge favours whichever answer it sees first), \textbf{verbosity bias} (longer answers score higher independent of quality), and \textbf{self-enhancement bias} (a model judges its own outputs more favourably). Han et al.'s \textquotedblleft{}Judge's Verdict: A Comprehensive Analysis of LLM Judge Capability Through Human Agreement\textquotedblright{} (arXiv:2510.09738, 2025) goes further, arguing a judge should not be trusted on its raw score at all — it should first be benchmarked against a measured human-agreement rate, the way a new diagnostic test is validated against a gold standard before deployment. \textbf{Relevance:} this is the literature-level grounding for a decision already made implicitly in Section~\ref{sec:gap} (point 2) and Section~\ref{sec:m-harness} — that an LLM-judge component in my evaluation harness cannot be trusted at face value — and it means any future LLM-judge step added to \texttt{src/evaluation/} needs either a structural cross-check or its own measured agreement rate against InnoBeweegLab's human reviewer, not just a plausible-sounding prompt.



\textbf{Reproducibility is a methodology in itself, with known, specific pitfalls.} Biderman et al., \textquotedblleft{}Lessons from the Trenches on Reproducible Evaluation of Language Models\textquotedblright{} (arXiv:2405.14782, EleutherAI, authors of the widely-used \texttt{lm-eval} harness) document — from three years of running a shared evaluation framework across the field — that evaluation results are routinely non-reproducible because of unreported \textquotedblleft{}minor\textquotedblright{} implementation details (prompt formatting, answer-extraction rules), shifting conventions, and social incentives to under-report negative or inconvenient results. Their recommended fixes: release the evaluation code itself, document the exact methodology, report variance across multiple repeated runs rather than a single score, and run exploratory error analysis rather than just a headline number. \textbf{Relevance:} this is the literature precedent for practices my harness already follows by construction (\texttt{src/evaluation/} is itself the released evaluation code; \texttt{runner.py}/\texttt{results\_log.py} already log mean$\pm$stdev across repeats, per Section~\ref{sec:m-harness}) — the earlier draft of this review had adopted the practice without citing why it is good practice.



\textbf{Human evaluation has its own, separately-documented best-practice literature — and the field frequently fails to follow it.} Van der Lee et al.'s \textquotedblleft{}Best practices for the human evaluation of automatically generated text\textquotedblright{} (INLG 2019) is the foundational checklist for designing a human-rater study of generated text: pre-registering evaluation criteria, piloting the rating instrument, reporting rater demographics/training, and reporting inter-rater agreement. Mei et al.'s \textquotedblleft{}Illusions of the Gold Standard: A Large-scale Analysis of Human Evaluation Protocols for Long-form Text Generation\textquotedblright{} (arXiv:2606.07936, 2026) is a direct, large-scale check on whether the field actually follows that checklist: across 284 manually-annotated *CL papers (2023–2025) plus 1,800+ more, \textquotedblleft{}only around half... include guidelines for human evaluation tasks,\textquotedblright{} and \textbf{no paper in the sample reported a power analysis to justify its sample size.} \textbf{Relevance:} this is a direct warning for Section~\ref{sec:m-reference}'s human-rubric step — as originally drafted, that step did not state its own evaluation guidelines, rater qualification, or sample-size justification in writing, which is exactly the gap Mei et al. find in roughly half of the published literature. Section~\ref{sec:m-reference} is revised below to close this specific gap rather than repeat it.



\subsection{Field-phase evidence: my own InnoBeweegLab stakeholder survey (n=5) — results relevant to validating the research gap and sharpening the main research question}\label{sec:res-survey}

Everything in Sections~\ref{sec:res-generation}--\ref{sec:res-evalmethod} is literature or production-system evidence from outside InnoBeweegLab. This subsection is different: it is primary field data I collected myself, a \textquotedblleft{}Research \& Reporting Workflow Survey\textquotedblright{} distributed to InnoBeweegLab researchers, coordinators, and project leads via Microsoft Forms. As of 7 October 2026 it has 5 responses and remains open, so the numbers below are indicative at this sample size, not final — I will re-run this analysis with the full response set before the final report. Two individual questions (data types normally worked with; difficulty combining qual/quant/expert findings) returned data that looks like a Forms display quirk rather than a real signal (nearly every respondent selected \textquotedblleft{}Other\textquotedblright{} on a straightforward multiple-choice data-type question) and are flagged rather than reported as findings.



\textbf{Workload and pain points.} The analysis-and-reporting phase is a real time cost: 4 of 5 respondents spend 21+ hours per project on it (2 at 21–40h, 2 at more than 40h), and for 3 of 5 the writing/structuring step alone takes 11–20 hours. Combining information from different data sources takes 2–10 hours for most respondents. Ranked by how many respondents flagged each as a current problem (of 5, multi-select): \textquotedblleft{}Writing reports takes too much time\textquotedblright{} and \textquotedblleft{}Maintaining consistency across reports\textquotedblright{} tie at 4; \textquotedblleft{}Data is stored in different formats\textquotedblright{} and \textquotedblleft{}Reports require significant rewriting\textquotedblright{} tie at 3; \textquotedblleft{}Combining sources is difficult\textquotedblright{} and \textquotedblleft{}Analysis takes too much time\textquotedblright{} tie at 2; \textquotedblleft{}It is difficult to trace conclusions back to evidence\textquotedblright{} scored only 1 — the lowest item on the list.



\begin{calloutbox}{amber}{Design implication}

Traceability and human control: not today's complaint, but tomorrow's requirement. That last number is worth pausing on, because what respondents say they currently struggle with and what they say actually matters point in opposite directions. Asked directly how important it is that every AI-generated conclusion can be traced back to its supporting evidence, the average score was \textbf{4.8 out of 5}; asked how important it is that a researcher can review, edit, approve, or reject AI-generated content before it becomes part of a report, the average was a perfect \textbf{5.0 out of 5} — unanimous. Read together with the problem-ranking above, the finding is: nobody is struggling to trace sources today because nobody is really attempting AI-assisted drafting yet. Traceability and human control are not reported pain points, they are the non-negotiable condition for adopting a new tool at all — exactly the premise my main research question is built on, now with quantified stakeholder backing rather than a single anecdote (Section~\ref{sec:gap}).

\end{calloutbox}

\textbf{What respondents want — and don't want — from an AI-supported tool.} Four functions got unanimous (5/5) support: analysing Excel/quantitative data, summarising interviews or qualitative responses, identifying themes and patterns, and comparing findings across multiple data sources — exactly the synthesis-heavy steps the workload data above says cost the most hours. Interest drops for \textquotedblleft{}generate draft findings/conclusions\textquotedblright{} (3/5 each), drops further for \textquotedblleft{}generate draft recommendations\textquotedblright{} (1/5): respondents want help synthesising evidence, not help deciding what to recommend. \textquotedblleft{}Supporting review and editing of generated content\textquotedblright{} scored 0/5, which — read against the 5.0/5 importance score above — looks like respondents treat human review as an assumed baseline rather than a feature worth asking for, not as something nobody wants.



\textbf{What they're worried about — and what they're not.} Top concerns, by how many respondents flagged each: incorrect information, unsupported conclusions/hallucinations, and generic recommendations each scored 4/5, followed by important context being missed and bias at 3/5 each. Compliance-style concerns scored far lower: privacy/GDPR, data security, and lack of transparency each got only 1/5. Respondents are worried about the system being \textit{wrong}, not about it being \textit{unsafe to use} — a direct steer toward where the evaluation-methodology work in Sections~\ref{sec:res-evalmethod} and~\ref{sec:m-reference} should concentrate its rigour.



\textbf{A concrete success bar, in their own words.} Asked what reduction in analysis/reporting time would make such a system worthwhile, 3 of 5 said 25–50\% and 2 of 5 said 10–25\%; nobody asked for more than 50\%. That is a real, stakeholder-stated number I can fold into the \textquotedblleft{}most efficient, cost-effective, and highest-value\textquotedblright{} framing of my main research question. Free-text answers reinforce the same theme on both sides: what would build trust is seeing \textquotedblleft{}on which data the recommendation is based\textquotedblright{} and knowing \textquotedblleft{}the sources actually check out\textquotedblright{}; what would make the system unacceptable is \textquotedblleft{}mistakes in the conclusions,\textquotedblright{} or — as one respondent put it — if \textquotedblleft{}correcting/checking the AI data costs more time than writing it ourselves,\textquotedblright{} which sets a hard floor under the time-savings number above.



4 of 5 respondents are willing to participate in a follow-up interview, and all 5 are willing to test and evaluate a prototype later in the project — concrete recruitment for the Phase 4 stakeholder evaluation (Section~\ref{sec:m-reference}) this report's own evaluation methodology depends on.



\subsection{Maps, charts and infographics as evidence --- results relevant to \textit{analysing the map inputs the way the reference report does}}\label{sec:res-maps}

Every method in Sections~\ref{sec:res-generation}--\ref{sec:res-pipelines} grounds text in text. None of them cites a region of an image. Yet InnoBeweegLab's data perspective reaches the report as images (Section~\ref{sec:process-perspectives}):
\begin{itemize}
\item maps on which the answer to an RIVM rule is a colour per building, a distance zone, or a marker;
\item infographics with printed numbers;
\item a sunburst chart.
\end{itemize}
The user supplies the maps as finished images, so the tool has to read them. This subsection collects what the literature says about reading such images automatically, and how reliable that is.

\begin{itemize}
\item \textbf{Local vision-language models read printed text and numbers well.} The Qwen2.5-VL technical report (Bai et al., 2025) gives its 7B model 95.7 on DocVQA, 87.3 on ChartQA and 82.6 on InfoVQA. A 7B model is in the same size class as the drafting models considered for the 16GB-RAM laptop. ChartQA (Masry et al., 2022) is a benchmark of 9.6K human-written and 23.1K generated questions about charts. The same model reports bounding boxes and points in the coordinates of the input image, which would let a reading point to the image region it came from. \textbf{Relevance:} the numbers printed on InnoBeweegLab's demographic infographics are within reach of a laptop-class model.
\item \textbf{The same models fail at simple spatial judgements.} \textquotedblleft{}Vision language models are blind\textquotedblright{} (Rahmanzadehgervi, Bolton, Taesiri \& Nguyen, ACCV 2024) tested four frontier models (GPT-4o, Gemini-1.5 Pro, Claude-3 Sonnet, Claude-3.5 Sonnet) on visual tasks that are easy for people. Their average accuracy was:
\begin{itemize}
\item 56.84\% on whether two lines intersect;
\item 39.44\% on counting overlapping circles;
\item 48.90\% on following single-coloured paths;
\item 58.07\% across all tasks, against 24\% for random guessing.
\end{itemize}
The authors trace the failures to how images are tokenised, and name checking whether streets intersect on a map as a real-world consequence. \textbf{Relevance:} \textquotedblleft{}does this building lie inside the 200\,m zone?\textquotedblright{} is exactly this kind of judgement.
\item \textbf{Map reasoning is an open problem.} MapBench (Xing et al., 2025) poses 1,649 path-finding queries on 100 maps of nine types. Under both zero-shot and chain-of-thought prompting, it \textquotedblleft{}significantly challenges\textquotedblright{} Llama-3.2-11B-Vision, Qwen2-VL-7B, GPT-4o and GPT-4o mini. MapEval (Dihan et al., ICML 2025) asks 700 questions about maps of 180 cities in 54 countries and tests 30 foundation models. No model exceeds 67\% accuracy, all trail humans by more than 20 percentage points, and they are weakest on distances, directions and routes.
\item \textbf{Extracting first and reasoning second works far better than asking the model directly.} MapQA (Chang, Palzer, Li, Fosler-Lussier \& Xiao, NeurIPS 2022 TRL workshop) builds a dataset of 62,367 choropleth maps with 795,525 questions. Its V-MODE method first extracts the map's underlying data, the value of every region, into a table, and only then answers questions over that table. For maps with a discrete legend the extraction reaches 92.0--99.9\% accuracy. Question answering on top of the table beats end-to-end baselines: on the MapQA-S subset (with oracle text recognition) it scores 87.1, against 72.2 for the best baseline and 89.8 for humans. DePlot (Liu et al., Findings of ACL 2023) applies the same idea to charts: it first translates a plot into a table, then lets a language model reason over it. This improves on the fine-tuned state of the art by 29.4\% on human-written ChartQA questions.
\item \textbf{Repeated readings can be turned into an uncertainty signal.} Self-ensembling (Berkane, Wang \& Majumder, 2026) prompts the same vision-language model several times on one chart, aligns the resulting tables, and takes the median of every cell. On a chart-extraction benchmark this improves accuracy by up to 8.09 points (about 23\% relative). The spread between readings correlates with error (Spearman $\rho$ between $-0.30$ and $-0.37$), which the authors present as a heuristic uncertainty measure, not a calibrated probability.
\item \textbf{Working on the GIS data instead of the picture.} Autonomous GIS (Li \& Ning, 2023) uses GPT-4 as a reasoning engine that runs GIS operations on spatial data (data collection, analysis, mapping) and produced correct results in three case studies. This is the stronger route when vector data is available. It does not fit this project, because it depends on a hosted model and InnoBeweegLab's input to the reporting step is the finished map image.
\end{itemize}

\begin{calloutbox}{teal}{Relevance for InnoBeweegLab}

The decisive information on InnoBeweegLab's maps is encoded in colour per building, in distance zones and in markers (Table~\ref{tab:rules-maps}). Those are the fine spatial judgements vision-language models get wrong. The infographics, by contrast, carry printed numbers, which these models read well. The literature therefore points to a split. The maps should be measured deterministically by reading the legend into a table, as in MapQA and DePlot. Model reading should be kept for printed numbers, with repeated readings and a human check (Section~\ref{sec:res-hitl}). Section~\ref{sec:m-maps} describes how I apply this.

\end{calloutbox}
